\section{Draft standard operating procedures}\label{app:sops}
These drafts are a starting point. The student turns them into the group's SOPs during WP0--WP3, with
the actual equipment, part numbers and parameters.

\subsection*{SOP-1: End-face inspection and cleaning}
\begin{enumerate}
  \item Wear gloves; work on a clean mat; have caps ready.
  \item Inspect the end face with the inspection microscope at the standard settings. Save the image as
        \texttt{<fibreID>\_<end>\_<YYYYMMDD>\_<step>.png}.
  \item If clean (no defects in the core zone, nothing loose in the cladding zone): stop and cap.
  \item Dry clean (filtered gas, then a dry cleaning cassette or click cleaner); inspect and image.
  \item If still contaminated: wet--dry clean with reagent-grade isopropanol on a lint-free wipe, then
        immediately a dry wipe; inspect and image. Repeat at most twice.
  \item Record every step in the fibre's log. Never clean a fibre that is scheduled for surface analysis.
\end{enumerate}

\subsection*{SOP-2: Polishing and acceptance of a test end}
\begin{enumerate}
  \item Polish with the group's procedure (films, sequence, times, pressure, lubricant written down) and
        log all parameters.
  \item Clean and inspect (SOP-1). Measure end-face geometry if an interferometer is available.
  \item Measure the coupling efficiency at low power (\SI{\le 0.3}{mW}) with the standard launch.
  \item Accept if the core zone is defect-free and the efficiency is within the agreed band around the
        batch median; otherwise re-polish. Cap and store the fibre in a clean container until use.
\end{enumerate}

\subsection*{SOP-3: Cell qualification (leak and oxygen hold test)}
\begin{enumerate}
  \item Clean all parts (isopropanol wipe, no silicone grease), bake printed parts under \Ntwo{} purge,
        assemble with a dummy connector.
  \item Pressure decay: pressurise to \SI{+100}{mbar} with \Ntwo, close the valves, log pressure and
        temperature for \SI{12}{h}. Correct for temperature ($\Delta p/p=\Delta T/T$) and compute
        $Q=V\,\mathrm{d}p/\mathrm{d}t$. Accept if $Q<\SI{1e-5}{mbar.L/s}$.
  \item Oxygen hold (static cells): purge with \Ntwo{} to $\lesssim\SI{1}{\percent}$ \Oxy, insert the
        absorber, seal, and log the \Oxy{} sensor spot reading for at least \SI{7}{days}. Accept if \Oxy{}
        stays below \SI{100}{ppm} after the initial absorption period.
  \item Record the cell ID, material, print settings, sealing and test results in the cell log.
\end{enumerate}

\subsection*{SOP-4: Start of a run}
\begin{enumerate}
  \item Fibres accepted (SOP-2), IDs and ends assigned to arms by randomisation, assignment saved in the
        run metadata.
  \item Cells qualified (SOP-3); purge flows set and logged; \Oxy, RH and VOC readings stable.
  \item Detectors zeroed; reference power meter cross-check recorded.
  \item Couple each arm at low power, then raise to the test power and set equal powers; record the
        initial values.
  \item Start the logger; write the run plan (duration, re-peak schedule, stop criteria) into the notebook.
\end{enumerate}

\subsection*{SOP-5: End of a run}
\begin{enumerate}
  \item Record the final throughput before and after re-peaking at the test power, then at low power.
  \item Switch off the laser. Inspect both ends without touching (SOP-1, imaging only).
  \item Either run the cleaning sequence of \cref{sec:ladder} (output end first) with a low-power
        measurement after each step, or cap the fibre and store it in an \Ntwo-filled, labelled container
        for surface analysis.
  \item Archive the data, images and notebook entries of the run and write a one-page run summary.
\end{enumerate}

\section{Data-logging specification}\label{app:logging}
\begin{table}[h]
\centering\small
\caption{Minimum channel list for the logger. Raw sampling at \SI{1}{Hz}; stored as 1-minute means
with standard deviations; one file per day plus a metadata file per run.}
\label{tab:channels}
\begin{tabularx}{\textwidth}{@{}l Y l@{}}
\toprule
Channel & Description & Units \\
\midrule
\texttt{t\_utc} & timestamp (ISO 8601, UTC) & -- \\
\texttt{arm<k>\_ref}, \texttt{arm<k>\_out}, \texttt{arm<k>\_back} & photodiode voltages of arm $k$ & V \\
\texttt{arm<k>\_O2}, \texttt{arm<k>\_RH}, \texttt{arm<k>\_VOC} & atmosphere of arm $k$ (cell or outlet) & ppm, \%, index \\
\texttt{arm<k>\_flow}, \texttt{arm<k>\_T} & purge flow, temperature at the cell & sccm, \si{\celsius} \\
\texttt{laser\_I}, \texttt{laser\_T}, \texttt{laser\_P} & diode current, temperature, monitor & mA, \si{\celsius}, V \\
\texttt{lab\_T}, \texttt{lab\_RH}, \texttt{lab\_p} & lab environment & \si{\celsius}, \%, hPa \\
\texttt{event} & free-text event log entries (re-peak, swap, cleaning, outage) & -- \\
\bottomrule
\end{tabularx}
\end{table}
The run metadata (JSON or YAML) contains: run ID, start and end time, fibre IDs and ends per arm,
atmosphere per arm, gas cylinder IDs and certificates, cell IDs, test power, detector calibration
factors, software version (git commit) and the analysis plan version.
